\chapter{绪论}

\section{研究背景与意义}

近年来，数字化和智能化进程对社会生活、经济发展等方面产生了巨大影响，教育领域也正经历着深刻的变革。特别是以大规模开放在线课程（MOOC）为代表的在线教育的快速发展，推动了优质教育资源和功能服务的跨时空共享，有效促进了教育公平，为学习者的均衡发展和持续学习创造了条件和机会\cite{whu-bachelor:1}。然而，随着在线教育平台中学习资源的不断累积，也暴露出了一些典型的痛点和难点问题：一方面，现有的平台在课程导航和个性化推荐方面略显不足，学习者在海量资源面前往往感到无所适从，难以快速定位符合自身需求的学习内容；另一方面，教育活动本身具有多维性和复杂性，学习者的背景知识、学习目标和认知能力存在显著差异，但多数平台仍采用“一刀切”的标准课程模式，无法为不同的学习者提供动态的、个性化的学习路径\cite{whu-bachelor:2}。此外，现有的代码实训和学习效果评估模块多采用简单的用例比对，只能给出冰冷的“对/错”判定，缺乏如人类教师般循循善诱的启发式指导，学习者很难从错误中获得深层次的认知提升。

针对上述问题，构建一个基于大语言模型（LLM）\cite{whu-bachelor:7,whu-bachelor:9,whu-bachelor:11}辅助的个性化学习系统显得尤为重要。个性化学习（Personalized Learning）强调以学习者为中心，根据其知识储备、兴趣偏好和学习目标，制定独一无二的学习策略，从而提高学习效率和效果。本课题旨在设计并实现一个大模型辅助个性化学习系统，通过引入协同过滤算法\cite{whu-bachelor:6}、检索增强生成（RAG）技术\cite{whu-bachelor:8}以及基于大模型的代码启发式分析算法\cite{whu-bachelor:10}，打破传统在线教育平台的局限，为学习者提供智能化的导学、深度的代码实训反馈以及个性化的课程内容推荐，这不仅在理论上具有重要的研究价值，在教育实践中也具有广阔的应用前景。

\section{国内外研究现状}

\subsection{竞品分析：同类产品现状}
MOOC的兴起在推动教育创新方面起到了重要作用。国外的Coursera、edX和Udacity三大平台凭借其丰富的课程资源和多元的教学模式吸引了全球数千万的学习者。Coursera通过穿插测验提供即时反馈；edX通过丰富的互动社区增强了学习体验；Udacity则通过项目驱动的课程设计将视频与练习紧密结合\cite{whu-bachelor:3}。国内方面，“学堂在线”和“中国大学MOOC”等平台也取得了显著发展，提供了大量本土化的优质课程。然而，这些主流平台在个性化学习路径定制和智能启发式辅导方面仍处于探索阶段，多数评测系统尚未实现深度语义理解和启发式反馈。

\subsection{关键技术分析：同类关键技术}
个性化学习的核心在于准确获取学习者的知识状态并进行有效建模\cite{whu-bachelor:4}。传统的方法如经典测试理论（CTT）和项目反应理论（IRT）主要从宏观角度评估学习者的综合能力。近年来，随着机器学习的发展，认知诊断（CD）和知识追踪（KT）等方法被广泛应用于微观知识状态的评估。深度知识追踪模型（DKT）\cite{whu-bachelor:5}和图结构知识追踪模型（GKT）进一步提高了状态预测的准确性。在智能问答和知识检索方面，多数平台仍使用传统的倒排索引检索（如BM25\cite{whu-bachelor:14}），面对多模态数据和复杂的自然语言提问时，语义理解能力较差。引入检索增强生成（RAG）结合大语言模型已成为一种新的技术趋势，它能结合FAISS等高效的向量数据库\cite{whu-bachelor:13}以及BGE等稠密向量模型\cite{whu-bachelor:12}，提供准确的知识回溯。在代码分析方面，目前的OJ系统多依赖黑盒测试，而基于AST的白盒静态分析技术结合大模型推理\cite{whu-bachelor:15}能够从语义层面上发现代码的逻辑漏洞和规范性问题，为学习者提供更深度的反馈。

\subsection{开发工具与平台分析}
目前大部分在线教育平台的架构主要基于传统的微服务或单体架构，前端多采用React或Vue等主流框架，后端则采用Java Spring Boot或Python Django等。本系统通过采用前后端分离的架构，前端使用Vue 3和Vite，结合Pinia状态管理和Element Plus组件库，提供了高性能和良好的用户体验。后端则采用Python的Flask框架，这得益于Python在人工智能和机器学习领域的生态优势，能够更加便捷地集成LangChain等大模型中间件、基于AST的代码分析模块以及各类异步任务处理队列。通过这样的技术选型和架构设计，本系统不仅具有良好的扩展性和维护性，更凸显了在智能化教育应用场景中的先进性与特色。

\section{本文的研究内容}

本文主要针对现有在线学习平台的不足，设计并实现一个大模型辅助的个性化学习系统，主要研究内容包括：

（1）\textbf{系统需求分析与架构设计}：分析个性化学习平台在不同用户角色（学生、教师）下的功能需求和非功能需求，设计基于Vue 3和Flask的前后端分离系统架构\cite{whu-bachelor:24,whu-bachelor:25}。

（2）\textbf{智能问答与RAG技术集成}：实现基于大语言模型和混合检索（BM25与语义向量检索）的智能问答助手，支持对课程视频和文档的高效知识检索，并通过前端溯源组件提供精准的知识定位\cite{whu-bachelor:27}。

（3）\textbf{启发式智能代码实训模块的实现}：设计并实现基于静态抽象语法树（AST）分析与LLM推理的代码评估算法。系统针对学习者的代码提交，不直接提供正确答案，而是基于多轮提交限制提供苏格拉底式的反问与启发，培养学习者的调试与重构能力\cite{whu-bachelor:19}。

（4）\textbf{个性化推荐系统开发}：实现基于协同过滤算法的课程推荐功能，结合学习者的历史行为、学习路径和打分记录，挖掘潜在的兴趣课程\cite{whu-bachelor:21}。

\section{本文的组织结构}

本文共分为七章，具体安排如下：
第一章为引言，介绍研究背景、意义、国内外研究现状及本文的研究内容；
第二章为国内外现状，深入分析竞品、同类关键技术及开发工具与平台；
第三章为系统需求分析，从功能性和非功能性两个维度详细描述系统需求，并给出用例规约；
第四章为系统设计，涵盖系统架构设计、核心模块设计与数据库设计；
第五章为系统实现，重点阐述启发式智能代码实训算法及RAG混合检索算法的实现细节，字数不少于 3000 字且包含算法描述；
第六章为系统测试，展示系统在功能、性能、白盒、黑盒等方面的测试结果及界面展示；
第七章为总结与展望，总结本文工作并对未来的研究方向提出展望。
